\chapter{Appendix A}
\label{app:appendix_a}

This appendix compiles supplementary technical material to support reproducibility and methodological transparency for the experiments detailed in Chapters~\ref{ch:baseline}, \ref{ch:nonlinear}, and \ref{ch:gw_alignment}. Section~\ref{sec:app_environment} documents the complete software and hardware environment utilized. Section~\ref{sec:app_script_mapping} maps the thirteen-test ablation study of Chapter~\ref{ch:gw_alignment} to the specific scripts and execution modes that produced each result. Section~\ref{sec:app_pseudocode} presents formal pseudocode for the core Optimal Transport algorithms previously described only in prose. Section~\ref{sec:app_hyperparam_spaces} consolidates the comprehensive hyperparameter search spaces explored during the Chapter~\ref{ch:gw_alignment} ablation study. Section~\ref{sec:app_notation} provides a reference table of the mathematical notation employed throughout this thesis. Finally, Section~\ref{sec:app_data_availability} concludes with a summary of data and code availability.

\section{Software and Hardware Environment}
\label{sec:app_environment}

All experiments in this thesis were executed on the University of Padova's \gls{hpc} cluster, \texttt{labsrv7}, utilizing Slurm for job scheduling and resource allocation. Table~\ref{tab:app_hw} details the compute hardware utilized across the various stages of the pipeline, while Table~\ref{tab:app_sw} specifies the software environment and package versions underlying all reported results.

\begin{table}[htpb]
\centering
\small
\begin{tabularx}{\textwidth}{lX}
\toprule
\textbf{Resource} & \textbf{Specification} \\
\midrule
Submit / access host & \texttt{labsrv7.math.unipd.it} \\
GPU nodes available & \texttt{dellcuda1} (1$\times$ NVIDIA A100), \texttt{dellcuda0}, \texttt{dellcuda2}, \texttt{dellsrv1} (1$\times$ NVIDIA V100 each), \texttt{dellsrv2--dellsrv5} (1$\times$ NVIDIA T4 each) \\
CPU-only nodes (lightweight verification jobs) & e.g.\ \texttt{hpblade04--hpblade16} (Test 12--13 sweeps) \\
Job scheduler & Slurm workload manager \\
\bottomrule
\end{tabularx}
\caption{Hardware resources used across the experimental pipeline, per
the Department of Mathematics HPC cluster specification
(\url{https://hpc.math.unipd.it/cluster-description/}). Jobs were
submitted from the \texttt{labsrv7} access host and scheduled onto the
cluster's compute nodes via Slurm; CPU allocation for Tests 12--13 is
confirmed directly from job logs (\texttt{Devices: [CpuDevice(id=0)]}).}
\label{tab:app_hw}
\end{table}

\begin{table}[htpb]
\centering
\small
\begin{tabularx}{\textwidth}{lX}
\toprule
\textbf{Component} & \textbf{Version} \\
\midrule
Python                          & 3.10 \\
Conda environment               & \texttt{ott\_env} \\
JAX                             & 0.6.2 \\
OTT-JAX                         & 0.6.0 \\
POT (Python Optimal Transport)  & 0.9.6 \\
Pymanopt                        & (RISGW Stiefel manifold optimisation) \\
scikit-learn                    & (StandardScaler, CCA cross-validation utilities) \\
SciPy                           & (SVD, linear-sum-assignment for Hungarian matching) \\
\bottomrule
\end{tabularx}
\caption{Software environment (\texttt{ott\_env} conda environment) under
which all experiments in Chapters~\ref{ch:baseline}--\ref{ch:gw_alignment}
were executed.}
\label{tab:app_sw}
\end{table}

Feature files are stored on the cluster filesystem within the \texttt{features\_100k/} directory as individual \texttt{.npy} arrays per encoder (\texttt{feat\_\{encoder\}\_part\_*.npy}). These were pre-partitioned into batches to facilitate the auto-resume mechanism described in Section~\ref{subsec:feature_pipeline}. All stochastic operations --- including anchor/query splitting, synthetic data generation, and projection-direction sampling for SGW --- were explicitly seeded via \texttt{numpy.random.RandomState}. The specific seed values are reported alongside their corresponding results throughout Chapters~\ref{ch:baseline}--\ref{ch:gw_alignment}.

\section{Script-to-Experiment Mapping}
\label{sec:app_script_mapping}

The thirteen-test ablation study detailed in Section~\ref{sec:ablation_study} and the core GW variant implementations of Section~\ref{sec:gw_methods} were executed using a consolidated set of dedicated Python scripts. Each script supports multiple execution modes selectable via command-line arguments. Table~\ref{tab:app_script_mapping} provides the complete mapping between every quantitative result reported in Chapter~\ref{ch:gw_alignment} and the specific script and mode utilized to generate it, ensuring exact reproducibility.

\begin{table}[htpb]
\centering
\small
\resizebox{\textwidth}{!}{%
\begin{tabular}{lll}
\toprule
\textbf{Script} & \textbf{Mode / Function} & \textbf{Thesis Result} \\
\midrule
\texttt{test\_lrgw\_exhaustive.py} & \texttt{--mode sanity}      & Test 1 (Table~\ref{tab:sanity_check}) \\
\texttt{test\_lrgw\_exhaustive.py} & \texttt{--mode selfgw}      & Test 2 (Table~\ref{tab:self_gw}) \\
\texttt{test\_lrgw\_exhaustive.py} & \texttt{--mode eps\_sweep}  & Test 3 (\S\ref{subsec:hyperparameter_sweeps}) \\
\texttt{test\_lrgw\_exhaustive.py} & \texttt{--mode sgw\_vs\_lrgw} & Test 4, Test 10 (Table~\ref{tab:cross_library}) \\
\texttt{test\_lrgw\_exhaustive.py} & \texttt{--mode raw\_gw}     & Test 5, Test 12 (Table~\ref{tab:raw_features_test}) \\
\texttt{test\_lrgw\_exhaustive.py} & \texttt{--mode cca\_sweep}  & Test 6, Test 8 (\S\ref{subsec:dataset_generalization}) \\
\texttt{test\_lrgw\_exhaustive.py} & \texttt{--mode all\_pairs}  & Test 9 (\S\ref{subsec:dataset_generalization}) \\
\texttt{test\_lrgw\_exhaustive.py} & \texttt{--mode transport}   & Test 11 (\S\ref{subsec:library_analysis}) \\
\texttt{test\_lrgw\_exhaustive.py} & \texttt{--mode eps\_low}    & Test 13 (\S\ref{subsec:advanced_verifications}) \\
\texttt{test\_sgw\_exhaustive.py}  & grid search (SGW/RISGW/SS-RISGW) & \S\ref{subsec:sliced_gw} best configurations \\
\texttt{test\_fgw\_after\_cca.py}  & FGW, FGW anchor-constrained & \S\ref{subsec:fused_gw} (Table~\ref{tab:gw_results}) \\
\texttt{gw\_diagnostic.py}         & six-pair compatibility sweep & Table~\ref{tab:gw-diagnostic} \\
\bottomrule
\end{tabular}%
}
\caption{Complete mapping from thesis results (Chapter~\ref{ch:gw_alignment})
to the generating script and execution mode. Feature paths, hyperparameters,
and seeds for each mode are as specified in the corresponding chapter
section.}
\label{tab:app_script_mapping}
\end{table}

Each script outputs its numerical results to a corresponding \texttt{results\_*.json} file (e.g., \texttt{results\_lrgw\_exhaustive.json}, \texttt{results\_gw\_diagnostic.json}). These files were subsequently parsed to populate the tables and figures presented throughout this thesis.

\section{Algorithm Pseudocode}
\label{sec:app_pseudocode}

This section presents formal pseudocode for the two core algorithms underlying the unsupervised alignment methods discussed in Chapter~\ref{ch:gw_alignment}: the Sliced Gromov-Wasserstein approximation from Section~\ref{subsec:sliced_gw}, and the Sinkhorn matrix-scaling procedure from Section~\ref{subsec:sinkhorn}, which forms the foundation of the LRSinkhorn sub-solver utilized by the Low-Rank Gromov-Wasserstein method (Section~\ref{subsec:low_rank_gw}).

\begin{algorithm}[htpb]
\caption{Sliced Gromov-Wasserstein (SGW) Score Accumulation}
\label{alg:sgw}
\begin{algorithmic}[1]
\Require 3D features $X \in \mathbb{R}^{n \times k}$, text features $Y \in \mathbb{R}^{n \times k}$, number of projections $P$
\Ensure Score matrix $S \in \mathbb{R}^{n \times n}$
\State $S \gets \mathbf{0}_{n \times n}$
\For{$p = 1$ \textbf{to} $P$}
    \State Sample random unit vector $\theta \sim \mathrm{Unif}(S^{k-1})$
    \State $x_\theta \gets X \theta$ \Comment{project onto $\theta$}
    \State $y_\theta \gets Y \theta$
    \State $\pi_X \gets \mathrm{argsort}(x_\theta)$ \Comment{1D ordering}
    \State $\pi_Y \gets \mathrm{argsort}(y_\theta)$
    \State $c_X \gets$ pairwise squared distances of $\mathrm{sort}(x_\theta)$
    \State $c_Y \gets$ pairwise squared distances of $\mathrm{sort}(y_\theta)$
    \State $\mathrm{cost}_{\mathrm{id}} \gets \| c_X - c_Y \|^2$
    \State $\mathrm{cost}_{\mathrm{anti}} \gets \| c_X - \mathrm{reverse}(c_Y) \|^2$
    \If{$\mathrm{cost}_{\mathrm{id}} \leq \mathrm{cost}_{\mathrm{anti}}$}
        \State $S[\pi_X, \pi_Y] \gets S[\pi_X, \pi_Y] + 1$
    \Else
        \State $S[\pi_X, \mathrm{reverse}(\pi_Y)] \gets S[\pi_X, \mathrm{reverse}(\pi_Y)] + 1$
    \EndIf
\EndFor
\State \Return $S$
\end{algorithmic}
\end{algorithm}

Algorithm~\ref{alg:sgw} accumulates votes over $P$ random 1D projections based on the optimal 1D ordering implied by each projection --- selecting either the identity alignment of the two sorted sequences or their reversal, depending on which yields a lower structural cost. The resulting score matrix $S$ functions directly as a similarity matrix for Top-$K$ retrieval and Hungarian matching evaluation, consistent with the methodology in Section~\ref{subsec:matching_acc}. This approach closely follows Algorithm 1 proposed by \textcite{vayer2019optimal}, as referenced in Section~\ref{subsec:sliced_gw}.

\begin{algorithm}[htpb]
\caption{Sinkhorn Matrix Scaling (entropic OT sub-solver)}
\label{alg:sinkhorn}
\begin{algorithmic}[1]
\Require Cost matrix $C \in \mathbb{R}^{n \times n}$, marginals $\mathbf{a}, \mathbf{b} \in \mathbb{R}^n$, regularisation $\epsilon > 0$, max.\ iterations $I$
\Ensure Transport plan $T \in \mathbb{R}^{n \times n}$
\State $K \gets \exp(-C / \epsilon)$ \Comment{Gibbs kernel}
\State $u \gets \mathbf{1}_n$, $v \gets \mathbf{1}_n$
\For{$i = 1$ \textbf{to} $I$}
    \State $u \gets \mathbf{a} \oslash (K v)$ \Comment{element-wise division}
    \State $v \gets \mathbf{b} \oslash (K^\top u)$
    \If{$\| u \odot (Kv) - \mathbf{a} \|_1 < \text{tol}$}
        \State \textbf{break} \Comment{marginal constraints satisfied}
    \EndIf
\EndFor
\State $T \gets \mathrm{diag}(u) \, K \, \mathrm{diag}(v)$
\State \Return $T$
\end{algorithmic}
\end{algorithm}

Algorithm~\ref{alg:sinkhorn} is the linear sub-solver called at every outer iteration of the Low-Rank Gromov-Wasserstein solver (Section~\ref{subsec:low_rank_gw}). Because the quadratic $\mathcal{GW}$ objective is non-convex, OTT-JAX's \texttt{GromovWasserstein} solver handles it by iteratively linearising the problem into a sequence of ordinary entropic \gls{ot} sub-problems of exactly this form, each one solved via a series of Sinkhorn iterations before the outer loop moves the linearisation point. The regularisation parameter $\epsilon$ here is the identical hyperparameter swept in Tests 3 and 13 (Sections~\ref{subsec:hyperparameter_sweeps} and~\ref{subsec:advanced_verifications}); the \texttt{LRSinkhorn} variant used throughout Chapter~\ref{ch:gw_alignment} additionally restricts $u$ and $v$ to act on a rank-$r$ factorisation of $K$ rather than the full $n \times n$ matrix, as described in Section~\ref{subsec:low_rank_gw}.

\section{Full Hyperparameter Search Spaces}
\label{sec:app_hyperparam_spaces}

While Chapter~\ref{ch:gw_alignment} reports the optimal configuration for each hyperparameter sweep within the thirteen-test ablation study, the comprehensive search spaces are distributed throughout the text. This section consolidates every hyperparameter range evaluated in Chapter~\ref{ch:gw_alignment} into three organized reference tables, categorized by their function: Table~\ref{tab:app_search_core} covers the core optimisation hyperparameters that apply broadly across LR-GW regardless of variant --- rank, regularisation, subspace dimension, and query size, swept in Tests 3, 6, 7, and 8. Table~\ref{tab:app_search_methods} reports the hyperparameters specific to the sliced and semi-supervised GW variants (FGW, SS-RISGW, SGW, RISGW), including the large-scale grid search underlying the SGW family. Table~\ref{tab:app_search_preprocessing} covers the two preprocessing choices tested in Tests 4 and 12. Unless stated otherwise, ``best value'' always refers to the configuration achieving the highest Top-5 retrieval accuracy on PointNet$\times$CLIP.

\begin{table}[htpb]
\centering
\small
\begin{tabularx}{\textwidth}{p{3.0cm}Xp{4.2cm}l}
\toprule
\textbf{Hyperparameter} & \textbf{Full Search Space} & \textbf{Best Value} & \textbf{Source} \\
\midrule
LR-GW rank $r$ & $\{2, 5, 10, 20, 50, 100, 200\}$ & $r{=}50$ (2.4\% Top-5) & Test 6 \\
Entropic reg.\ $\epsilon$ (standard) & $\{0.001, \ldots, 5.0\}$, 8 values & identical cost for all & Test 3 \\
Entropic reg.\ $\epsilon$ (fine, low range) & $\{10^{-4}, \ldots, 10^{-1}\}$, 7 values & identical cost for all & Test 13 \\
CCA subspace dimension $k$ & $\{10, 20, 50, 100\}$ & $k{=}20$ for LR-GW (3.0\% Top-5); $k{=}50$ main protocol & Test 8 \\
Query set size $n$ & $\{500, 1000, 2000\}$ & $n{=}500$ (no gain from more data) & Test 7 \\
\bottomrule
\end{tabularx}
\caption{Core optimisation hyperparameters swept during the ablation study
of Chapter~\ref{ch:gw_alignment} (Tests 3, 6, 7, 8), independent of the
specific GW variant used.}
\label{tab:app_search_core}
\end{table}

\begin{table}[htpb]
\centering
\small
\resizebox{\textwidth}{!}{%
\begin{tabular}{llll}
\toprule
\textbf{Method} & \textbf{Hyperparameter} & \textbf{Search Space / Best Value} & \textbf{Source} \\
\midrule
FGW & Trade-off $\alpha$ & $[0, 1]$ continuous; best $\alpha{=}0.2$ (4.4\% Top-5) & \S\ref{subsec:fused_gw} \\
FGW anchor-constrained & Penalty & $\{0, 10, 100\}$; best $=10$ (identical to unconstrained) & \S\ref{subsec:fused_gw} \\
SS-RISGW & Anchor weight $\lambda$ & $\{0, 0.1, 1.0, 10.0\}$; best $\lambda{=}1.0$ (1.4\% Top-5) & \S\ref{subsec:sliced_gw} \\
SGW & Projection count & 100 (fixed) & \S\ref{subsec:sliced_gw} \\
RISGW & Projections / iterations & 20 proj., 30 grad.\ iter.\ (fixed) & \S\ref{subsec:sliced_gw} \\
\midrule
\multicolumn{4}{l}{\textit{SGW / RISGW / SS-RISGW exhaustive grid search (288 configurations total)}} \\
\midrule
\multicolumn{4}{l}{$k \in \{20, 50, 100\}$, \quad $n_{\text{anchor}} \in \{1\text{k}, 5\text{k}, 30\text{k}\}$, \quad $\lambda \in \{0, 0.1, 1.0, 10.0\}$, \quad 3 seeds} \\
\multicolumn{4}{l}{Best: $k{=}20$, $n{=}5{,}000$ (SGW, 3.2\% Top-5) --- see \texttt{test\_sgw\_exhaustive.py}} \\
\bottomrule
\end{tabular}%
}
\caption{Method-specific hyperparameters and their search spaces for the
semi-supervised and sliced GW variants (FGW, SS-RISGW, SGW, RISGW),
including the large-scale 288-configuration grid search.}
\label{tab:app_search_methods}
\end{table}

\begin{table}[htpb]
\centering
\small
\begin{tabularx}{\textwidth}{Xp{4.0cm}Xl}
\toprule
\textbf{Choice} & \textbf{Search Space} & \textbf{Best Value} & \textbf{Source} \\
\midrule
Cost function & $\{$squared-Euclidean, angular (arccos)$\}$ & angular (1.9\% vs.\ 0.4\% Top-5) & Test 4 \\
Feature normalisation (pre-CCA) & $\{L_2\text{-norm}, \text{StandardScaler}\} \times \{\text{no SM}, \text{SM}\}$ & all $\approx$ equivalent ($\approx$1.5\%) & Test 12 \\
\bottomrule
\end{tabularx}
\caption{Preprocessing choices tested in Tests 4 and 12. ``SM'' denotes
the median-ratio scale-matching normalisation introduced by
\textcite{li2026scale} (Section~\ref{subsec:unsupervised_alignment}).}
\label{tab:app_search_preprocessing}
\end{table}

\section{Notation and Symbols}
\label{sec:app_notation}

Table~\ref{tab:app_notation} summarizes the principal mathematical notation employed throughout this thesis to facilitate cross-chapter reference.

\begingroup
\small
\begin{longtable}{lll}
\caption{Principal mathematical notation used throughout this thesis, in order of first appearance.}
\label{tab:app_notation} \\
\toprule
\textbf{Symbol} & \textbf{Meaning} & \textbf{First used} \\
\midrule
\endfirsthead

% Intestazione per le pagine successive
\multicolumn{3}{c}%
{{\tablename\ \thetable{} -- continued from previous page}} \\
\toprule
\textbf{Symbol} & \textbf{Meaning} & \textbf{First used} \\
\midrule
\endhead

% Piè di pagina per le pagine spezzate
\midrule
\multicolumn{3}{r}{{\small\textit{Continued on next page}}} \\
\endfoot

% Piè di pagina finale
\bottomrule
\endlastfoot

% --- Contenuto della tabella ---
$X, Y$                  & 3D and text feature matrices & \S\ref{subsec:supervised_alignment} \\
$X^r, Y^r$               & CCA-projected 3D and text features & \S\ref{subsec:supervised_alignment} \\
$W_X, W_Y$               & CCA projection matrices & \S\ref{subsec:supervised_alignment} \\
$k$                      & CCA subspace dimensionality & \S\ref{subsec:supervised_alignment} \\
$R, b$                   & Affine transformation weight and bias & \S\ref{subsec:supervised_alignment} \\
$n_A$                    & Number of anchor pairs & \S\ref{subsec:anchor_protocol} \\
$n$                      & Number of query samples (or GW sample size) & \S\ref{sec:metrics} \\
$\mu, \nu$                & Source and target probability measures & \S\ref{subsec:monge_kantorovich} \\
$\pi, T$                 & Optimal transport / transport plan (coupling matrix) & \S\ref{subsec:monge_kantorovich} \\
$c(x,y), C$               & Transport cost function / cost matrix & \S\ref{subsec:monge_kantorovich} \\
$W_p$                    & $p$-Wasserstein distance & \S\ref{subsec:monge_kantorovich} \\
$\epsilon$                & Entropic regularisation strength & \S\ref{subsec:sinkhorn} \\
$H(T)$                   & Entropy of a transport plan & \S\ref{subsec:sinkhorn} \\
$u, v$                   & Sinkhorn scaling vectors & \S\ref{subsec:sinkhorn} \\
$\mathcal{GW}(X,Y)$        & Gromov-Wasserstein distance & \S\ref{subsec:gw_definition} \\
$d_X, d_Y$                & Intra-domain distance metrics & \S\ref{subsec:gw_definition} \\
$L(\cdot,\cdot)$          & GW discrepancy loss function & \S\ref{subsec:gw_definition} \\
$r$                      & Low-rank GW factorisation rank & \S\ref{subsec:low_rank_gw} \\
$Q, R, g$                 & LR-GW low-rank factor matrices & \S\ref{subsec:low_rank_gw} \\
$\theta$                 & Random projection direction (SGW) & \S\ref{subsec:sliced_gw} \\
$\Delta$                 & RISGW rotation matrix (Stiefel manifold) & \S\ref{subsec:sliced_gw} \\
$\lambda$                & SS-RISGW anchor penalty weight & \S\ref{subsec:sliced_gw} \\
$\alpha$                 & FGW structural/feature trade-off & \S\ref{subsec:fused_gw} \\
$s$                      & Median-ratio scale-matching factor & \S\ref{subsec:unsupervised_alignment} \\
$M, L$                    & RML Mahalanobis matrix and Cholesky factor & \S\ref{subsec:rml_implementation} \\
$K, \gamma$                & KCCA kernel matrix and RBF bandwidth & \S\ref{subsec:kcca_implementation} \\
$m$                      & Nystr\"{o}m landmark count & \S\ref{subsec:kcca_implementation} \\
$\rho(T)$                 & Diagonal-to-uniform concentration ratio & \S\ref{subsec:entropy_definition} \\
$\mathrm{CKA}(K,L)$        & Centered Kernel Alignment score & \S\ref{subsec:cka_definition} \\
\end{longtable}
\endgroup

\section{Data and Code Availability}
\label{sec:app_data_availability}

All scripts underlying the feature extraction pipeline, alignment method implementations, and diagnostic ablation study are stored on the \texttt{labsrv7} cluster filesystem, organised within a centralised project directory rather than an external code repository. Raw extracted features for all six encoder combinations (Table~\ref{tab:app_sw}) are retained there within the \texttt{features\_100k/} directory. These features are not redistributed as part of this thesis due to their substantial storage requirements ($100{,}000$ objects across up to $1280$ dimensions per encoder). The Objaverse dataset \cite{deitke2023objaverse} and its corresponding Cap3D captions \cite{luo2023scalable} are publicly accessible from their original sources. Similarly, all pre-trained encoder weights utilized in this research (PointNet, OpenCLIP ViT-bigG-14, BERT, RoBERTa) have been released by their respective authors. The sole exception is SparseConv, for which pre-trained weights were unavailable at the time of writing (Section~\ref{subsec:3d_encoders}). The numerical results underlying the tables in Chapters~\ref{ch:baseline}--\ref{ch:gw_alignment} are preserved in the corresponding \texttt{results\_*.json} output files, as referenced in Table~\ref{tab:app_script_mapping}.